%=============================================================
%  第4章 結果
%=============================================================
\chapter{結果}
\label{chap:results}

\section{クラスターの特性}
\label{sec:results-clusters}

K-Means クラスタリングにより，日本上場企業は表\ref{tab:cluster-profile}に示す
5つのクラスターに分類された。

\begin{table}[htbp]
  \centering
  \caption{クラスター財務プロファイル（中央値，2010〜2024年度）}
  \label{tab:cluster-profile}
  % TODO: \input{tables/tab_cluster_profile.tex} に置き換える
  \small
  \begin{tabular}{llrrrrr}
    \toprule
    & & \multicolumn{5}{c}{クラスター} \\
    \cmidrule(lr){3-7}
    変数 & 単位 & C0 & C1 & C2 & C3 & C4 \\
    \midrule
    総資産     & 億円 & 208  & 51   & 616  & 3,867 & 55 \\
    自己資本比率 & \%  & 42.2 & 43.9 & 68.7 & 42.9  & 67.2 \\
    構成比率   & \%   & 32.5 & 9.9  & 18.0 & 21.3  & 18.3 \\
    \midrule
    解釈 & & 中小規模 & 小規模 & 中規模 & 大企業型 & 小規模 \\
         & & 標準型  & 脆弱型 & 優良型 &         & 健全型 \\
    \bottomrule
  \end{tabular}
  \begin{tablenotes}
    \small
    \item 注: 構成比率は全観測値に占める各クラスターの割合。
  \end{tablenotes}
\end{table}

各クラスターの特徴を以下に述べる。

\paragraph{\Cluster{0}（中小規模・標準型）}
総資産の中央値は約 208 億円であり，自己資本比率 42.2\%，営業利益は正の値を示す。
全観測値の 32.5\% を占める最大クラスターであり，日本上場企業の標準的な財務構造を体現する。

\paragraph{\Cluster{1}（小規模・財務脆弱型）}
総資産の中央値は約 51 億円と最小規模であり，自己資本比率は 43.9\% である。
PC1・PC2 ともに全クラスター中で最低値を示し，特に営業利益が負である傾向が顕著である。
全観測値の 9.9\% を占める。

\paragraph{\Cluster{2}（中規模・財務優良型）}
自己資本比率が 68.7\% と高く，財務の健全性が際立つ。

\paragraph{\Cluster{3}（大企業型）}
総資産の中央値は約 3,867 億円と突出して大きく，大企業グループを構成する。

\paragraph{\Cluster{4}（小規模・財務健全型）}
規模は小さいが，自己資本比率 67.2\% と高い財務健全性を持つ。

\section{クラスター遷移行列}
\label{sec:results-transition}

表\ref{tab:transition-matrix}に年度間のクラスター遷移行列を示す。

\begin{table}[htbp]
  \centering
  \caption{クラスター遷移行列（行: 遷移元，列: 遷移先, \%）}
  \label{tab:transition-matrix}
  % TODO: \input{tables/tab_transition_matrix.tex} に置き換える
  \begin{tabular}{lrrrrr}
    \toprule
    & \multicolumn{5}{c}{遷移先クラスター} \\
    \cmidrule(lr){2-6}
    遷移元 & C0 & C1 & C2 & C3 & C4 \\
    \midrule
    C0 & \textbf{79.9} & 5.2  & 6.0  & 4.1  & 4.8 \\
    C1 & 22.0 & \textbf{61.3} & 1.9  & 0.9  & 13.9 \\
    C2 & 9.0  & 1.1  & \textbf{80.6} & 4.1  & 5.2 \\
    C3 & 4.9  & 0.6  & 4.1  & \textbf{90.4} & 0.0 \\
    C4 & 9.4  & 7.2  & 5.7  & 0.0  & \textbf{77.7} \\
    \bottomrule
  \end{tabular}
  \begin{tablenotes}
    \small
    \item 注: 各行の合計は 100\%。太字は対角成分（自己遷移）。
    \item 全サンプルの自己遷移率は 80.1\%。
  \end{tablenotes}
\end{table}

対角成分（自己遷移率）は全クラスター平均で 80.1\% であり，
年度をまたいでも財務クラスターの帰属が高い安定性を持つことが示される。

\section{遷移要因の分析}
\label{sec:results-probit}

\subsection{二値プロビット}

\begin{table}[htbp]
  \centering
  \caption{二値プロビット推定結果}
  \label{tab:probit}
  % TODO: \input{tables/tab_probit_results.tex} に置き換える
  \begin{tabular}{lrr}
    \toprule
    変数 & モデル1（年度 FE なし） & モデル2（年度 FE あり） \\
    \midrule
    PC1         & — & $-0.008$ \\
    PC2         & — & $-0.105^{***}$ \\
    PC3         & — & $0.022^{**}$ \\
    \midrule
    Pseudo $R^2$ & — & 0.043 \\
    N           & — & \num{46084} \\
    \bottomrule
  \end{tabular}
  \begin{tablenotes}
    \small
    \item 注: $^{*}p<0.1$, $^{**}p<0.05$, $^{***}p<0.01$
  \end{tablenotes}
\end{table}

\section{予測モデルの結果}
\label{sec:results-prediction}

\subsection{遷移の二値分類（Step 1）}

表\ref{tab:binary-performance}に，遷移の有無を予測するモデルの性能を示す。

\begin{table}[htbp]
  \centering
  \caption{Binary 分類モデルの性能（Test: 2022〜2024 年度）}
  \label{tab:binary-performance}
  \begin{tabular}{lrrl}
    \toprule
    モデル & Accuracy & AUC-ROC & 備考 \\
    \midrule
    Logistic Regression & 0.62 & 0.630 & ベースライン \\
    XGBoost             & 0.69 & \textbf{0.757} & \texttt{scale\_pos\_weight=3.99} \\
    \bottomrule
  \end{tabular}
\end{table}

XGBoost は AUC = 0.757 を達成し，ロジスティック回帰（AUC = 0.630）を大幅に上回った。

\subsection{SHAP 分析: 遷移要因}

図\ref{fig:shap-binary-importance}に，遷移の二値分類に対する特徴量重要度を示す。

\begin{figure}[htbp]
  \centering
  \includegraphics[width=0.70\linewidth]{shap_binary_importance.png}
  \caption{特徴量重要度（Mean $|\mathrm{SHAP}|$）：遷移有無（Binary XGBoost）}
  \label{fig:shap-binary-importance}
\end{figure}

最も重要な特徴量は営業利益（Mean $|\mathrm{SHAP}|$ = 0.311）であり，
次いで長期借入金（0.297），現在クラスター（0.269）が続く。

\subsection{遷移先の多クラス分類（Step 2）}

遷移した企業のみを対象とする多クラス分類の結果を表\ref{tab:mc-performance}に示す。

\begin{table}[htbp]
  \centering
  \caption{Multi-class 分類モデルの性能（Test: 2022〜2024 年度，遷移企業のみ）}
  \label{tab:mc-performance}
  \begin{tabular}{lrrl}
    \toprule
    モデル & Accuracy & AUC-ROC（macro） & 備考 \\
    \midrule
    XGBoost & 0.73 & \textbf{0.947} & 遷移企業のみ（$n=2,081$）\\
    \bottomrule
  \end{tabular}
\end{table}

多クラス分類において AUC（macro） = 0.947 という高い予測精度が達成された。

\subsection{SHAP 分析: 遷移先決定要因}

図\ref{fig:shap-mc-heatmap}に，各遷移先クラスターへの SHAP 値のヒートマップを示す。

\begin{figure}[htbp]
  \centering
  \includegraphics[width=\linewidth]{shap_multiclass_heatmap.png}
  \caption{Mean $|\mathrm{SHAP}|$ ヒートマップ：遷移先クラスター × 特徴量}
  \label{fig:shap-mc-heatmap}
\end{figure}

\section{ロバストネスチェックの結果}
\label{sec:results-robustness}

本節では，PCA の推定期間（全期間 vs 訓練期間のみ）と現在クラスター変数の有無を
組み合わせた 4 パターンの比較分析，および統計的有意性の検定結果を報告する。

\subsection{4 パターン比較}

表\ref{tab:robustness-4pattern}に，PCA 推定期間とクラスター変数の有無を
組み合わせた 4 パターンの予測性能を示す。
図\ref{fig:robustness-4pattern}にその比較を可視化する。

\begin{table}[htbp]
  \centering
  \caption{ロバストネスチェック: 4 パターンの予測性能比較（AUC-ROC）}
  \label{tab:robustness-4pattern}
  \begin{tabular}{llrr}
    \toprule
    パターン & 設定 & Binary & Multi-class \\
    \midrule
    A（基準）& 全期間 PCA + クラスター変数あり & \textbf{0.757} & \textbf{0.947} \\
    B        & Train-only PCA + クラスター変数あり & 0.752 & 0.929 \\
    C        & 全期間 PCA + クラスター変数なし & 0.745 & 0.879 \\
    D        & Train-only PCA + クラスター変数なし & 0.745 & 0.879 \\
    \bottomrule
  \end{tabular}
  \begin{tablenotes}
    \small
    \item 注: Pattern A と B のクラスターラベルは，ハンガリアン法による最適リラベリング後の一致率が 98.4\% であり，実質的に同一のクラスター構造を捉えている。
  \end{tablenotes}
\end{table}

\begin{figure}[htbp]
  \centering
  \includegraphics[width=0.85\linewidth]{robustness_4pattern_comparison.png}
  \caption{4 パターンの AUC-ROC 比較（Binary・Multi-class）}
  \label{fig:robustness-4pattern}
\end{figure}

4 パターンの結果から，以下の知見が得られた。
第一に，PCA 推定期間の影響（A vs B）は Binary で 0.005，Multi-class で 0.018 と小さく，
情報漏洩の影響は限定的である。
第二に，現在クラスター変数を除外した場合（A vs C），Binary では 0.012 の低下にとどまり，
財務変数のみでも遷移の有無を相当程度予測できる。
一方，Multi-class では 0.068 の低下が見られ，現在クラスターが遷移先の決定に
重要な役割を果たすこと，すなわちパス依存性の存在を示唆する。
第三に，クラスター変数を除外した場合（C vs D），PCA 推定期間の差異は消失し，
情報漏洩の懸念が現在クラスター変数に限定されることが確認された。

\subsection{統計的有意性の検定}

表\ref{tab:robustness-stat}に，各パターン間の AUC 差に対する
DeLong 検定（Binary）および Bootstrap 95\% 信頼区間（2,000 回リサンプリング）の結果を示す。
図\ref{fig:robustness-stat}に検定結果を可視化する。

\begin{table}[htbp]
  \centering
  \caption{パターン間の AUC 差に対する統計的検定結果}
  \label{tab:robustness-stat}
  \small
  \begin{tabular}{llrrll}
    \toprule
    比較 & 検証内容 & $\Delta$AUC & $p$値 & Bootstrap 95\% CI & 判定 \\
    \midrule
    \multicolumn{6}{l}{\textit{Binary}} \\
    A vs B & PCA 推定期間 & $+0.005$ & 0.121 & $[+0.002, +0.008]$ & n.s. \\
    A vs C & クラスター変数（全期間）& $+0.012$ & 0.040 & $[+0.007, +0.017]$ & $*$ \\
    B vs D & クラスター変数（Train-only）& $+0.007$ & 0.198 & $[+0.003, +0.012]$ & n.s. \\
    C vs D & PCA 推定期間（クラスターなし）& $+0.000$ & 1.000 & $[+0.000, +0.000]$ & n.s. \\
    \midrule
    \multicolumn{6}{l}{\textit{Multi-class（Bootstrap のみ）}} \\
    A vs B & PCA 推定期間 & $+0.018$ & — & $[+0.015, +0.022]$ & $*$ \\
    A vs C & クラスター変数（全期間）& $+0.069$ & — & $[+0.063, +0.074]$ & $*$ \\
    B vs D & クラスター変数（Train-only）& $+0.050$ & — & $[+0.045, +0.056]$ & $*$ \\
    C vs D & PCA 推定期間（クラスターなし）& $+0.000$ & — & $[+0.000, +0.000]$ & n.s. \\
    \bottomrule
  \end{tabular}
  \begin{tablenotes}
    \small
    \item 注: $*$: 95\% 信頼区間が 0 を含まない。n.s.: 非有意。
    DeLong 検定は Binary のみ適用。Multi-class は Bootstrap CI で判定。
  \end{tablenotes}
\end{table}

\begin{figure}[htbp]
  \centering
  \includegraphics[width=0.85\linewidth]{robustness_stat_test.png}
  \caption{パターン間 AUC 差の統計的検定結果（DeLong $p$値・Bootstrap 95\% CI）}
  \label{fig:robustness-stat}
\end{figure}

DeLong 検定の結果，Binary における PCA 推定期間の影響（A vs B）は統計的に非有意（$p = 0.121$）であり，
情報漏洩が予測精度に実質的な影響を与えていないことが確認された。
クラスター変数の寄与（A vs C）は $p = 0.040$ で有意であるが，
Train-only PCA における同比較（B vs D）では非有意（$p = 0.198$）となり，
クラスター変数の寄与は PCA 推定期間に依存する可能性が示唆された。
Multi-class においては，Bootstrap CI に基づき PCA 推定期間の影響は有意であるものの，
絶対差は 0.018 にとどまり，予測モデルの頑健性が確認された。
